% ═══════════════════════════════════════════════════════════════
% ARBEITSBLATT: Verbos en -ar (Regelmäßige -ar Verben)
% ═══════════════════════════════════════════════════════════════
% Sequenz 4: Mi colegio | Block c: Verbos en -ar
% Gesamtpunktzahl: 20 Punkte
% ═══════════════════════════════════════════════════════════════

\documentclass[11pt, a4paper]{article}

% SW-MODUS TOGGLE
\newif\ifswmodus
\swmodusfalse

% PAKETE
\usepackage[utf8]{inputenc}
\usepackage[T1]{fontenc}
\usepackage[ngerman]{babel}
\usepackage{helvet}
\renewcommand{\familydefault}{\sfdefault}

\usepackage[margin=2cm]{geometry}
\usepackage{enumitem}
\usepackage{tabularx}
\usepackage{booktabs}
\usepackage{multirow}
\usepackage{array}

\usepackage{xcolor}
\usepackage{framed}
\usepackage{tcolorbox}
\tcbuselibrary{skins}

\usepackage{fancyhdr}
\usepackage{lastpage}
\usepackage{needspace}

% FARBEN
\definecolor{spanisch-color}{HTML}{c41e3a}
\definecolor{grau-color}{HTML}{888888}
\definecolor{hellgrau-color}{HTML}{f5f5f5}
\definecolor{orange-color}{HTML}{b35c00}

\definecolor{sw-dark}{HTML}{000000}
\definecolor{sw-gray}{HTML}{666666}
\definecolor{sw-white}{HTML}{ffffff}

\ifswmodus
    \colorlet{spanisch}{sw-dark}
    \colorlet{grau}{sw-gray}
    \colorlet{hellgrau}{sw-white}
    \colorlet{orange}{sw-dark}
\else
    \colorlet{spanisch}{spanisch-color}
    \colorlet{grau}{grau-color}
    \colorlet{hellgrau}{hellgrau-color}
    \colorlet{orange}{orange-color}
\fi

\newcommand{\fachfarbe}{spanisch}

% VARIABLEN
\newcommand{\thema}{Verbos en -ar}
\newcommand{\sequenz}{04}
\newcommand{\block}{c}
\newcommand{\fach}{Spanisch}
\newcommand{\klasse}{7}
\newcommand{\maxpunkte}{20}
\newcommand{\datum}{\today}

% KOPF- UND FUSSZEILE
\pagestyle{fancy}
\fancyhf{}
\renewcommand{\headrulewidth}{0.4pt}
\renewcommand{\footrulewidth}{0.4pt}

\lhead{\fach{} -- Klasse \klasse}
\chead{AB-\sequenz\block}
\rhead{\datum}

\lfoot{}
\cfoot{}
\rfoot{Seite \thepage\ von \pageref{LastPage}}

% BEFEHLE
\newcommand{\punktebox}[1]{\hfill\fbox{\small \_/#1}}
\newcommand{\luecke}[1]{\rule{#1}{0.4pt}}

\newtcolorbox{merkkasten}[1][]{
    colback=hellgrau,
    colframe=\fachfarbe,
    colbacktitle=white,
    coltitle=\fachfarbe,
    boxrule=1pt,
    arc=0pt,
    left=8pt, right=8pt, top=8pt, bottom=8pt,
    fonttitle=\bfseries,
    attach boxed title to top left={yshift=-2mm, xshift=5mm},
    boxed title style={boxrule=0pt, colback=white},
    title=#1
}

\newtcolorbox{vokabelkasten}{
    colback=white,
    colframe=\fachfarbe,
    boxrule=0.5pt,
    arc=0pt,
    left=5pt, right=5pt, top=5pt, bottom=5pt
}

\newtcolorbox{hinweis}{
    colback=white,
    colframe=orange,
    colbacktitle=white,
    coltitle=orange,
    boxrule=1pt,
    arc=0pt,
    left=5pt, right=5pt, top=5pt, bottom=5pt,
    fonttitle=\bfseries,
    attach boxed title to top left={yshift=-2mm, xshift=5mm},
    boxed title style={boxrule=0pt, colback=white},
    title=Hinweis
}

\newenvironment{aufgabe}[1]{%
    \needspace{5\baselineskip}%
    \nopagebreak[4]%
    \vspace{10pt}
    \noindent\textbf{\textcolor{\fachfarbe}{Aufgabe #1}}
    \nopagebreak[4]%
    \vspace{5pt}
    \nopagebreak[4]%
    \begin{list}{}{\leftmargin=0pt}
    \item[]
}{%
    \end{list}
}

\newlist{teilaufgaben}{enumerate}{2}
\setlist[teilaufgaben,1]{label=\alph*), leftmargin=2em}
\setlist[teilaufgaben,2]{label=\roman*), leftmargin=2em}

\newcommand{\es}[1]{\textcolor{\fachfarbe}{\textbf{#1}}}

% ═══════════════════════════════════════════════════════════════
% DOKUMENT
% ═══════════════════════════════════════════════════════════════

\begin{document}

% KOPFBEREICH
\begin{center}
    {\Large\bfseries\textcolor{\fachfarbe}{Arbeitsblatt: \thema}}
\end{center}

\vspace{5pt}

\noindent
\begin{tabularx}{\textwidth}{|X|X|X|}
\hline
\textbf{Name:} \luecke{4cm} &
\textbf{Klasse:} \klasse &
\textbf{Datum:} \luecke{2cm} \\
\hline
\end{tabularx}

\vspace{15pt}

% ═══════════════════════════════════════════════════════════════
% MERKKASTEN: Konjugation -ar Verben
% ═══════════════════════════════════════════════════════════════

\begin{merkkasten}[Merkkasten: Konjugation regelmäßiger -ar Verben]

\textbf{Bildung:} Infinitiv (z.B. \es{hablar}) $\rightarrow$ Stamm (\es{habl-}) + Endung

\vspace{0.5em}

\begin{center}
\begin{tabular}{|l|c|c|l|}
\hline
\textbf{Person} & \textbf{Endung} & \textbf{hablar} & \textbf{Deutsch} \\
\hline
yo & \luecke{1.5cm} & habl\luecke{1cm} & ich spreche \\[0.6em]
\hline
tú & \luecke{1.5cm} & habl\luecke{1cm} & du sprichst \\[0.6em]
\hline
él/ella/usted & \luecke{1.5cm} & habl\luecke{1cm} & er/sie spricht \\[0.6em]
\hline
nosotros/-as & \luecke{1.5cm} & habl\luecke{1cm} & wir sprechen \\[0.6em]
\hline
\end{tabular}
\end{center}

\vspace{0.5em}

\textbf{Wichtige -ar Verben:}

\vspace{0.3em}

\begin{tabular}{ll@{\hspace{2cm}}ll}
\es{hablar} & = \luecke{3cm} & \es{trabajar} & = \luecke{3cm} \\[0.6em]
\es{estudiar} & = \luecke{3cm} & \es{cantar} & = \luecke{3cm} \\[0.6em]
\es{escuchar} & = \luecke{3cm} & \es{bailar} & = \luecke{3cm} \\
\end{tabular}

\vspace{0.5em}

\textit{Tipp: Übertrage die Endungen und Übersetzungen aus dem Lernpfad!}
\end{merkkasten}

\vspace{10pt}

% ═══════════════════════════════════════════════════════════════
% AUFGABE 1: Endungen einsetzen (4P)
% ═══════════════════════════════════════════════════════════════

\begin{aufgabe}{1: Endungen einsetzen}
Setze die richtigen Endungen ein. \punktebox{4}
\end{aufgabe}

\begin{vokabelkasten}
\textit{Endungen:} -o | -as | -a | -amos
\end{vokabelkasten}

\vspace{0.5em}

\noindent
\begin{tabular}{rl@{\hspace{1.5cm}}rl}
a) & Yo habl\luecke{1.5cm} español. & c) & Nosotros estudi\luecke{1.5cm} matemáticas. \\[1em]
b) & Tú escuch\luecke{1.5cm} música. & d) & María trabaj\luecke{1.5cm} mucho. \\
\end{tabular}

\vspace{15pt}

% ═══════════════════════════════════════════════════════════════
% AUFGABE 2: Verben konjugieren (4P)
% ═══════════════════════════════════════════════════════════════

\begin{aufgabe}{2: Verben konjugieren}
Konjugiere die Verben in der richtigen Form. \punktebox{4}
\end{aufgabe}

\noindent
\begin{tabular}{rl}
a) & (cantar, yo) $\rightarrow$ Yo \luecke{3cm} en el coro. \\[1em]
b) & (bailar, tú) $\rightarrow$ ¿Tú \luecke{3cm} salsa? \\[1em]
c) & (estudiar, él) $\rightarrow$ Pedro \luecke{3cm} inglés. \\[1em]
d) & (trabajar, nosotros) $\rightarrow$ Nosotros \luecke{3cm} en grupo. \\
\end{tabular}

\vspace{15pt}

% ═══════════════════════════════════════════════════════════════
% AUFGABE 3: Sätze bilden (6P)
% ═══════════════════════════════════════════════════════════════

\begin{aufgabe}{3: Sätze bilden}
Bilde Sätze mit den angegebenen Elementen. \punktebox{6}
\end{aufgabe}

\begin{hinweis}
Achte auf die richtige Verbform und Satzstellung (Subjekt -- Verb -- Objekt).
\end{hinweis}

\vspace{0.5em}

\noindent
\begin{tabular}{rlp{9cm}}
a) & yo / hablar / español & \\[0.3em]
   & & \luecke{10cm} \\[1.2em]

b) & tú / escuchar / al profesor & \\[0.3em]
   & & \luecke{10cm} \\[1.2em]

c) & María / estudiar / historia & \\[0.3em]
   & & \luecke{10cm} \\[1.2em]

d) & nosotros / cantar / una canción & \\[0.3em]
   & & \luecke{10cm} \\[1.2em]

e) & Pablo / trabajar / en clase & \\[0.3em]
   & & \luecke{10cm} \\[1.2em]

f) & tú / bailar / muy bien & \\[0.3em]
   & & \luecke{10cm} \\
\end{tabular}

\vspace{15pt}

% ═══════════════════════════════════════════════════════════════
% AUFGABE 4: Schulalltag beschreiben (6P)
% ═══════════════════════════════════════════════════════════════

\begin{aufgabe}{4: Schulalltag beschreiben}
Schreibe mindestens 4 Sätze über deinen Schulalltag. Verwende verschiedene -ar Verben. \punktebox{6}
\end{aufgabe}

\begin{hinweis}
Mögliche Verben: hablar, estudiar, escuchar, trabajar, cantar, bailar\\
Beispiel: \textit{Yo estudio español en el colegio.}
\end{hinweis}

\vspace{0.5em}

\noindent
\textbf{Mi día en el colegio:}

\vspace{0.5em}

\noindent
1. \luecke{14cm}

\vspace{1em}

\noindent
2. \luecke{14cm}

\vspace{1em}

\noindent
3. \luecke{14cm}

\vspace{1em}

\noindent
4. \luecke{14cm}

\vspace{1em}

\noindent
(Zusatz:) 5. \luecke{14cm}

% ═══════════════════════════════════════════════════════════════
% PUNKTEÜBERSICHT
% ═══════════════════════════════════════════════════════════════

\vspace{20pt}
\hrule
\vspace{10pt}

\begin{flushright}
\textbf{Punkte:} \luecke{1cm} / \maxpunkte
\end{flushright}

\end{document}
